\documentclass[10pt,conference]{IEEEtran}
\IEEEoverridecommandlockouts

\usepackage{cite}
\usepackage{amsmath,amssymb,amsfonts}
\usepackage{graphicx}
\usepackage{textcomp}
\usepackage{xcolor}
\usepackage{url}

\def\BibTeX{{\rm B\kern-.05em{\sc i\kern-.025em b}\kern-.08em
    T\kern-.1667em\lower.7ex\hbox{E}\kern-.125emX}}

\begin{document}

\title{Imagining Quantum Futures Through Everyday Life Stories}

\author{\IEEEauthorblockN{1\textsuperscript{st} Takuma Sano}
\IEEEauthorblockA{\textit{Graduate School of Media Design}\\
\textit{Keio University}\\
Yokohama, Japan\\
takuma.sano@kmd.keio.ac.jp}
\and
\IEEEauthorblockN{2\textsuperscript{nd} Reiko Iwasaki}
\IEEEauthorblockA{\textit{Graduate School of Media Design}\\
\textit{Keio University}\\
Yokohama, Japan\\
r.iwasaki@kmd.keio.ac.jp}
\and
\IEEEauthorblockN{3\textsuperscript{rd} Haruka Shiina}
\IEEEauthorblockA{\textit{Graduate School of Media Design}\\
\textit{Keio University}\\
Yokohama, Japan\\
shiha@kmd.keio.ac.jp}
\and
\IEEEauthorblockN{4\textsuperscript{th} Shota Nagayama}
\IEEEauthorblockA{\textit{Graduate School of Media Design}\\
\textit{Keio University}\\
Yokohama, Japan\\
shota@qitf.org}
}

\maketitle

\section*{1. BoF Title}

Imagining Quantum Futures Through Everyday Life Stories

\section*{2. BoF Abstract}

Quantum technologies have often been imagined as forces that may reshape future life and society. Policy goals, research programs, and technology roadmaps already suggest possible futures in which quantum technologies contribute to medicine, energy, mobility, security, optimization, and other domains. At the same time, some use cases may remain undiscovered if discussions rely only on existing policy scenarios, corporate roadmaps, or predefined application categories. This Birds of a Feather session invites participants to explore such possibilities through everyday life stories. Participants will be presented with short fictional scenes set in futures influenced by quantum technologies. Rather than immediately evaluating feasibility, they will imagine the gadgets, services, social systems, infrastructures, and surrounding practices that may exist in those worlds but are not explicitly described in the stories. The session aims to help participants express and share future visions that draw on their own expertise, interests, and professional experience. Through individual reflection, shared discussion, and synthesis, participants will consider not only what may exist in a quantum-influenced future, but also when such ideas may become meaningful and what technical or social conditions would be required to support them. The expected outcome is a set of shared future visions, use-case ideas, and discussion points that may inform later exploration of quantum applications, collaborations, and social implications.

\section*{3. BoF Authors}

The proposed BoF will be organized by the following members of the Graduate School of Media Design, Keio University, Yokohama, Japan.

\begin{itemize}
\item Takuma Sano, organizer and session moderator, takuma.sano@kmd.keio.ac.jp
\item Reiko Iwasaki, organizer, r.iwasaki@kmd.keio.ac.jp
\item Haruka Shiina, organizer, shiha@kmd.keio.ac.jp
\item Shota Nagayama, organizer and main contact, shota@qitf.org
\end{itemize}

Takuma Sano will serve as the session moderator and facilitate the opening framing, storytelling exercise, participant discussion, and final reflection. Shota Nagayama will serve as the main contact. The organizers will prepare the future stories, design the worksheets, support participant engagement, and summarize discussion themes after the session.

\section*{4. BoF Target Audience}

This BoF is intended for all Quantum Week participants who are interested in future societies involving quantum technologies. The target audience includes researchers and engineers working on quantum applications, people responsible for technology roadmaps, new business development, or social implementation in companies, and participants engaged in technology governance, policy, standardization, and ecosystem formation. It also welcomes experts and practitioners from application domains such as materials science, quantum chemistry, life sciences, medicine, machine learning, physics, and engineering. The session is especially suitable for participants who wish to think from their own expertise and experience about possible futures, rather than only from already defined technical roadmaps. No single disciplinary background is required. The BoF is designed to benefit from a mixture of technical, industrial, policy, design, and social perspectives, and welcomes participants who want to actively discuss how quantum technologies may become part of future everyday life, services, and social systems.

\section*{5. BoF Motivation and Objectives}

Future visions of quantum technologies have been articulated through policy, research programs, and industrial activities. For example, Japan's Quantum Future Society Vision presents goals such as economic growth, harmony between people and the environment, and richer everyday life, and discusses examples including compact high-sensitivity MRI, tailored medicine, renewable-energy allocation, MaaS, evacuation routing, and emergency logistics \cite{quantum_future_society_vision}. In industry, comprehensive future-society visions are not always stated explicitly, but alliance announcements and executive statements often suggest expectations that quantum technologies will be combined with AI, HPC, and existing research-and-development processes and embedded in real-world systems. These expectations interact with one another and shape a broader social imagination of quantum futures.

At the same time, as expectations grow, it becomes increasingly important to distinguish plausible future pathways from hype. DARPA's Quantum Benchmarking Initiative, for example, explicitly frames the need to separate hype from reality in assessing whether utility-scale quantum computers can be realized by 2033 \cite{darpa_qbi}. However, the existence of hype does not mean that imaginative futures should be dismissed. In past technological development, the act of imagining technologies and ways of life that do not yet exist has often helped create new research questions and social possibilities. Science Fiction Prototyping is one method for this purpose: science-based fiction can help people explore future technologies and their human and social implications \cite{johnson_sfp}, support dialogue through fictional characters and narrative scenarios \cite{burnam_fink}, and provide occasions for participants to explore possible futures \cite{draudt}.

The objective of this BoF is therefore to treat scientifically and technically grounded ``dream stories'' as shared objects for discussion. Participants will use future everyday-life scenes as starting points to express and share future visions and possibilities that they may not yet have verbalized. They will imagine products, services, and social systems that may exist around those scenes, discuss when such ideas might become relevant, and consider what should be started now, and with whom, to move toward them. The goal is not to validate feasibility during the session, but to create generative contact points among participants with different expertise and experiences, supporting future research, development, implementation, and collaboration around quantum applications.

\section*{6. BoF Relevance}

This BoF is relevant to IEEE Quantum Week because it provides a participatory setting for discussing future quantum applications beyond predefined sectors and technical benchmarks. It may offer new insights to participants involved in research, development, and social implementation of quantum applications, especially by allowing people from optimization, materials, chemistry, medicine, AI, service design, and other application areas to share future visions grounded in their own expertise and professional experience. Quantum Week brings together communities concerned with quantum computing, quantum engineering, quantum software, applications, and industry ecosystems. These communities need not only technical milestones, but also ways to discuss what kinds of services, products, infrastructures, and collaborations may emerge when quantum capabilities are embedded in broader social and technological systems.

The session also addresses a known translation challenge between academic research and industrial practice. In quantum software engineering, for instance, researchers have noted gaps between academic and industry goals and the need for collaboration around practical application problems \cite{zappin}. This BoF does not attempt to solve that gap directly, but it creates a space where participants can articulate possible future applications from multiple perspectives before they become formal requirements or business categories.

Many discussions of quantum applications begin with an already defined industrial problem or computational task and ask how quantum technologies might be applied. This BoF reverses part of that direction. Rather than focusing only on the potential capability of quantum technologies themselves, it asks what may appear around people's lives once quantum technologies, classical computing, digital services, physical products, services, and institutions are combined. For example, the potential of quantum-enabled decision support, materials search, chemical calculation, and drug discovery is widely recognized, but the everyday scenes and social systems that may emerge when such capabilities are embedded in products and services remain uncertain. By using future stories as common discussion objects, the session encourages participants to consider cross-layer systems, cross-domain services, and the forms of collaboration that may be necessary to make them meaningful.

\section*{7. BoF Proposal}

\subsection*{Guiding Question}

The guiding question of the session is: \emph{How can we imagine worlds that current quantum science and technology have not yet reached?} The BoF treats storytelling not as a prediction exercise, but as a method for opening a shared discussion space around future quantum applications.

\subsection*{Session Format}

The BoF is planned as a 90-minute interactive session.

\begin{itemize}
\item \textbf{Framing and audience engagement, 20 minutes.} The organizers will ground the discussion by introducing the current technological stage of quantum internet and distributed quantum computing, and by framing their diffusion as an ecosystem-building process involving research institutions, companies, infrastructure providers, application developers, and users. The organizers will also introduce exploratory work on use cases and problem settings for a distributed quantum-computing era. An interactive polling tool will ask participants how they have imagined futures involving quantum technologies, and the responses will be briefly visualized and discussed. The organizers will then introduce science fiction as one way to imagine worlds not yet reached by current science and technology, refer to examples of ``Quantum x something'' in speculative works, and explain how the workshop stories were created from lifestyle scenarios that emerged internally through participant activities in the organizers' previous project. Schumpeter's idea of ``new combinations'' will be used as a cue for combining existing application ideas, users, and everyday situations.
\item \textbf{Story-based individual exploration, 18 minutes.} Participants will first indicate application areas of interest through an interactive poll. The organizers will then introduce approximately five future stories. Each participant will receive a paper worksheet containing a fictional person, time, everyday scene, and a short story of about 100--150 words, with a wide free-writing area beside it. Participants will freely imagine products, services, gadgets, social systems, or surrounding situations that may exist near the scene, drawing on their own expertise, interests, and professional experience. They do not need to limit themselves to one domain, and may write multiple ideas across multiple fields. They may use words, sentences, diagrams, sketches, or arrows. At this stage, participants are not required to map their ideas immediately to a specific quantum technology; they first explore what may exist around the future everyday-life scene.
\item \textbf{Break, 10 minutes.}
\item \textbf{Sharing and discussion, 27 minutes.} Participants will share major ideas that emerged from their worksheets. When useful, the participant who generated an idea will explain its background, motivation, and connection to their field. The discussion will use two main prompts: first, \emph{when} is this idea about, what future time horizon does it imply, and what technical or social stage must be reached for it to make sense? Second, what should be done \emph{now}, and by whom, to support its realization? Participants will consider the roles of researchers, companies, policymakers, users, and other actors.
\item \textbf{Reflection and closing, 15 minutes.} The session will close with reflection on the work itself. Participants will discuss whether they imagined something they had not previously considered, whether they found something they would want to make or use, and what they noticed while thinking from future everyday-life stories. The organizers will return to the guiding question and reflect on whether participants were able to share possibilities not yet captured by roadmaps or existing use-case lists. The closing will also consider how the ideas and contact points generated in the BoF may connect to future quantum-application research, development, implementation, and collaboration.
\end{itemize}

\subsection*{Story Preparation}

The future stories used in the BoF will be created by the organizers. They will share two common future conditions. First, quantum computing is assumed to support some of the computation needed to address real-world problems, while being integrated with classical computing and existing systems. Second, some research outcomes from quantum-enabled materials search, chemical calculation, or related fields are assumed to be embedded in actual products or services.

The stories will also draw on ideas generated in prior ideation activities with students at a Japanese university. The organizers will identify recurring lifestyle changes across those ideas, such as people no longer being conscious of scheduling or waiting time, physical distance becoming less restrictive for work and life, or energy management receding from the foreground of everyday life. These recurring changes will be combined with fictional characters and everyday scenes to create future stories in which the influence of quantum technologies appears indirectly.

\begin{table}[t]
\caption{Examples of Idea-to-Story Translation}
\label{tab:stories}
\centering
\scriptsize
\setlength{\tabcolsep}{2pt}
\begin{tabular}{p{0.25\linewidth}p{0.28\linewidth}p{0.30\linewidth}}
\hline
Application ideas & Common lifestyle change & Future story direction \\
\hline
No scheduling concept; no delivery-time selection; no waiting at airports & People rarely notice coordination or waiting & Everyday coordination happens quietly behind the scenes \\
No business travel for meetings; less need to move to a large city; people can live anywhere & Physical distance less strongly limits work and life choices & People work and interact across distance with little awareness of distance itself \\
Devices that do not need charging; high-performance battery materials; local energy generation & Energy management becomes less visible in daily life & People are freed from managing the technical systems that support their gadgets and daily routines \\
\hline
\end{tabular}
\end{table}

The stories will be reviewed by experts in quantum information, quantum internet, and distributed quantum computing. The review will check the assumed technical premises, the relationship with classical computing and existing systems, and whether the stories contain excessive technological leaps or scientific contradictions.

\subsection*{Expected Outcomes}

The main outcome will be the sharing of experiences and future visions that participants may not have previously articulated explicitly. By encountering and discussing ideas from people with different roles and research backgrounds, participants may generate insights that would be difficult to produce within a single field. The session also invites participants to bring their fundamental interests, naive expectations, and personal motivations toward quantum technologies into conversation, creating new contact points for future research, development, implementation, and collaboration. The outputs will include participant-generated ideas, discussion notes, perceived time horizons, required technical and social conditions, and possible collaboration points. The session will not claim technical feasibility, deployment readiness, or quantum advantage. Instead, it will provide exploratory material that can later be examined through more formal readiness, impact, or application-evaluation methods.

\section*{8. BoF Biographies}

\textbf{Takuma Sano} is a master's student in the Graduate School of Media Design at Keio University and a member of Associate Professor Shota Nagayama's group. He holds a bachelor's degree in psychology from Rikkyo University, with a focus on Human--Machine Interaction research. He currently studies the mutual influence between people and their physical and informational environments, and investigates how quantum technologies can drive urban transformation.

\textbf{Reiko Iwasaki} is a master's student in the Graduate School of Media Design at Keio University and has served as a facilitator in previous workshops and contributed to projects that generate ideas about the future of quantum technologies. She also has experience in filmmaking, which informs her creative approach to visualizing future concepts.

\textbf{Haruka Shiina} received her B.A. in Cultural Anthropology from Tsuda University in 2025. In the same year, she entered the Graduate School of Media Design at Keio University and joined Q/est, the laboratory led by Associate Professor Shota Nagayama. Her research focuses on exploring how the social implementation of quantum technologies may transform human behavior and everyday life.

\textbf{Shota Nagayama} is an Associate Professor and the head of The Quantum Future Society Design Center at the Graduate School of Media Design, Keio University. He holds a Ph.D. in Media and Governance. He specializes in quantum internet, distributed quantum computing, and quantum error correction codes, particularly quantum internet/computer architectures and communication protocols. As a professional in the information engineering field, he has been collaborating with quantum physicists, founded and serves as the director of the Quantum Internet Task Force, and is the lead author of the white paper for QITF. He also serves as the Project Manager for a project demonstrating integrated quantum networks aimed at distributed quantum computing within the JST Moonshot R\&D Program, a large-scale national project aiming for Fault-Tolerant Quantum Computing in Japan, a board member of the WIDE project, and a secretary of the Information Processing Society of Japan's Special Interest Group on Quantum Software. His technical works have been awarded at IEEE Quantum Week twice.

\end{document}
